% ============================================================
%  DSM SF-PBL — STANDALONE COMPILE VERSION
%  Uses standard LaTeX article class — compiles immediately
%  without needing frontiersSCNS.cls
%
%  For final submission: copy content into official Frontiers
%  Word/LaTeX template from:
%  https://www.frontiersin.org/guidelines/author-guidelines
% ============================================================

\documentclass[12pt, a4paper]{article}

% ── Packages ─────────────────────────────────────────────────
\usepackage[utf8]{inputenc}
\usepackage[T1]{fontenc}
\usepackage[margin=2.5cm]{geometry}
\usepackage{setspace}
\usepackage{url}
\usepackage[hidelinks]{hyperref}
\usepackage{booktabs}
\usepackage{array}
\usepackage{multirow}
\usepackage{graphicx}
\usepackage{amsmath}
\usepackage{amssymb}
\usepackage{siunitx}
\usepackage{microtype}
\usepackage{natbib}
\usepackage{lineno}
\usepackage{xcolor}
\usepackage{parskip}
\usepackage{caption}
\usepackage{enumitem}
\usepackage{lmodern}
\usepackage{authblk}

\captionsetup{font=small, labelfont=bf}
\doublespacing
\linenumbers

% ── Title block ──────────────────────────────────────────────
\title{\textbf{A Computationally-Integrated SF-PBL Framework for
Engineering Dynamics: Design, Implementation, and Mixed-Methods
Evaluation in an OBE-Aligned Undergraduate Course}}

\author[1*]{Sunny Nanade}
\author[1]{Debasis Dash}
\author[1]{Sudipto Sarkar}
\author[1]{Koteswararao Anne}

\affil[1]{Mukesh Patel School of Technology Management \& Engineering, SVKM's NMIMS, Mumbai, India}

\affil[*]{Corresponding author: \href{mailto:sunny.nanade@nmims.edu}{sunny.nanade@nmims.edu} | ORCID: 0000-0001-7098-1084}

\date{Submitted to: \textit{Frontiers in Education} ---
STEM Education Section \\ \today}

% ─────────────────────────────────────────────────────────────
\begin{document}
\maketitle
\thispagestyle{empty}

% ── Running head ─────────────────────────────────────────────
\begin{center}
  \small\textit{Target Journal: Frontiers in Education |
  Article Type: Original Research | Word count (abstract): 247}
\end{center}

\vspace{1em}

% ============================================================
%  ABSTRACT
% ============================================================
\begin{abstract}
\noindent
Engineering dynamics courses demand simultaneous development of
mathematical modelling, computational implementation, and
communication competencies, yet conventional instruction rarely
provides authentic contexts for their integration. This study
presents the design, implementation, and mixed-methods evaluation
of a Simulation-Focused Problem-Based Learning (SF-PBL) framework
embedded within a 15-week undergraduate Dynamic System Modeling
(DSM) course ($N$\,=\,53, B.Tech Mechatronics Engineering, Semester
IV, January--April 2026) at an NBA-accredited institution. The
framework operationalises a four-phase pipeline --- System
Identification, Mathematical Modelling, Computational Simulation,
Validation \& Communication --- scaffolded through an
Explain-Try-Challenge model and assessed via a six-judge expert
panel. A pre--post mixed-methods design measured self-efficacy
(10-item Likert scale), conceptual knowledge (5-item MCQ), PBL
satisfaction, AI tool usage, and qualitative reflections.
Self-efficacy improved significantly (pre: $M$\,=\,2.92,
$SD$\,=\,0.21; post: $M$\,=\,3.80, $SD$\,=\,0.40;
$t$(52)\,=\,13.91, $p$\,$<$\,.001, Cohen's $d$\,=\,2.72, pooled
SD). Knowledge gains were large ($t$(52)\,=\,11.32,
$p$\,$<$\,.001, $d$\,=\,1.20; Hake's $g$\,=\,0.50, medium band).
All ten self-efficacy items improved significantly ($p$\,$<$\,.001).
PBL satisfaction was high ($M$\,=\,3.74, $SD$\,=\,0.30).
Eighty-five percent used AI tools, primarily for computational
coding. Qualitative themes identified simulation validation, PID
tuning, and linking derivation to computation as breakthrough
moments. Exhibition rubric scores from 17 student projects
($M$\,=\,82.62/100, $SD$\,=\,9.98) confirmed strong performance
across all domains. The SF-PBL framework provides a replicable,
evidence-based model for computationally-integrated dynamics
pedagogy in OBE-aligned STEM programmes.

\bigskip
\noindent\textbf{Keywords:} problem-based learning; engineering
dynamics; computational simulation; self-efficacy; OBE; authentic
assessment; AI-assisted learning; STEM education
\end{abstract}

\newpage

% ============================================================
%  1. INTRODUCTION
% ============================================================
\section{Introduction}

\subsection{The Pedagogical Challenge in Engineering Dynamics}

Engineering dynamics --- encompassing kinematics, constitutive
laws, Lagrangian mechanics, and vibration analysis --- occupies a
foundational position in mechanical and mechatronics engineering
curricula. It presents persistent pedagogical challenges: the
mathematical formalism is demanding, the physical intuition
difficult to develop, and the connection to real engineering
systems rarely made explicit within traditional lecture-driven
courses \citep{streveler2008learning}. Students frequently solve
textbook problems through surface-level pattern matching without
building the conceptual understanding required for engineering
practice.

Compounding this challenge is the increasing centrality of
computational tools. Modern engineers are expected not only to
derive equations of motion analytically but to implement them
numerically, simulate dynamic system behaviour, and interpret
computational outputs --- a set of competencies that conventional
closed-book assessments rarely evaluate authentically
\citep{justo2021enhancing}. The emergence of AI coding assistants
(ChatGPT, GitHub Copilot, Claude) has further transformed the
competency landscape: students can now generate simulation code
rapidly, raising important questions about depth of conceptual
understanding versus tool-assisted productivity.

\subsection{Problem-Based Learning as a Solution Pathway}

Problem-Based Learning (PBL), as defined by
\citet{hmelo2004problem}, places authentic, ill-structured problems
at the centre of the learning experience, requiring students to
self-direct their inquiry, build domain knowledge, and communicate
solutions collaboratively. In engineering education, PBL has been
shown to improve conceptual understanding, motivation, and
collaborative competencies \citep{prince2006inductive,
dym2005engineering}. However, its implementation in mathematically
intensive courses such as dynamics remains underexplored,
particularly when computational simulation is foregrounded as the
primary learning artefact.

\subsection{The Research Gap}

Three specific gaps motivate this work:

\begin{enumerate}
  \item \textbf{Computational-PBL integration gap}: Existing PBL
    implementations in engineering dynamics rarely position
    Python-based ODE simulation as the core learning product,
    leaving the derivation-to-computation bridge unaddressed.

  \item \textbf{Authentic assessment gap}: Multi-expert external
    evaluation panels covering industry, intellectual property,
    mathematics, and communication are uncommon in undergraduate
    PBL contexts; most studies rely on instructor-only assessment.

  \item \textbf{OBE alignment gap}: Few published studies
    demonstrate explicit, auditable traceability from Course
    Outcomes (COs) through Bloom's taxonomy levels to rubric
    dimensions within an accreditation-compliant (NBA/ABET)
    framework.
\end{enumerate}

\subsection{Purpose and Contributions}

This paper presents the design, implementation, and empirical
evaluation of the SF-PBL framework embedded within the Dynamic
System Modeling (DSM) course (702MH0C023), a 15-week core course
for B.Tech Mechatronics Engineering students at NMIMS MPSTME,
Mumbai (2 January -- 25 April 2026). The study makes four
contributions:

\begin{enumerate}
  \item A documented, replicable \textbf{four-phase SF-PBL
    computational algorithm} (System Identification
    $\rightarrow$ EOM Derivation $\rightarrow$ Python Simulation
    $\rightarrow$ Validation \& Communication)

  \item An \textbf{OBE-aligned rubric framework} with complete
    CO$\rightarrow$Bloom's$\rightarrow$NBA Programme Outcome
    traceability

  \item \textbf{Empirical evidence} of large-to-very-large
    learning gains (Cohen's $d$ $>$ 1.0 on both self-efficacy
    and knowledge) from a mixed-methods pre--post study
    ($N$\,=\,53)

  \item \textbf{Authentic assessment data} from a six-judge
    expert panel evaluating 17 student projects across five
    engineering domains
\end{enumerate}

% ============================================================
%  2. THEORETICAL FRAMEWORK
% ============================================================
\section{Theoretical Framework}

\subsection{Simulation-Focused Problem-Based Learning (SF-PBL)}

SF-PBL is defined in this study as a variant of PBL in which
\textbf{mathematical derivation from first principles and its
subsequent computational simulation constitute the primary learning
artefact}. This distinguishes SF-PBL from project-based learning
(where the product is typically a physical prototype) and from
standalone computational modelling courses (where system
identification and mathematical modelling steps are usually
pre-specified).

The SF-PBL framework rests on three theoretical pillars. First,
\textbf{Hmelo-Silver's PBL cycle} \citeyearpar{hmelo2004problem}:
students encounter an ill-structured real engineering system,
engage in self-directed inquiry to identify governing physics,
construct knowledge through mathematical and computational means,
and reflect collaboratively on outcomes. Second,
\textbf{Bandura's self-efficacy theory}
\citeyearpar{bandura1997self}: mastery experiences --- successfully
deriving equations of motion and validating them in simulation ---
are the most potent source of self-efficacy gains. The sequential
phases of SF-PBL are designed to provide escalating mastery
experiences mapped directly to course competencies. Third,
\textbf{Hake's interactive engagement framework}
\citeyearpar{hake1998interactive}: courses achieving normalised
gain $g$ $>$ 0.30 are classified as effectively
interactive-engagement; SF-PBL's four-phase structure, with active
derivation, implementation, and public defence, is hypothesised to
achieve medium-to-high normalised gain.

\subsection{OBE Framework and NBA Alignment}

The DSM course was designed using Outcome-Based Education (OBE)
principles \citep{spady1994outcome} with explicit backward design
from Programme Outcomes (POs) to Course Outcomes (COs) to
assessment rubric dimensions. Five Course Outcomes (CO1--CO5) are
detailed in Table~\ref{tab:cos}.

\begin{table}[ht]
\centering
\caption{Course Outcomes and Bloom's Taxonomy Levels}
\label{tab:cos}
\begin{tabular}{p{0.5cm} p{8cm} c}
\toprule
\textbf{CO} & \textbf{Statement} & \textbf{Bloom's Level} \\
\midrule
CO1 & Apply kinematic analysis to multi-body systems    & Apply (L3)    \\
CO2 & Construct FBDs for complex mechanisms             & Analyse (L4)  \\
CO3 & Derive EOMs using Newton's laws                   & Analyse (L4)  \\
CO4 & Apply Lagrangian mechanics to holonomic systems   & Evaluate (L5) \\
CO5 & Implement ODE solvers in Python and validate      & Create (L6)   \\
\bottomrule
\end{tabular}
\end{table}

\subsection{The Explain-Try-Challenge (E-T-C) Instructional Model}

Each 90-minute lecture session followed a three-phase model:
\textbf{Explain} (30 min) --- conceptual introduction with worked
examples; \textbf{Try} (40 min) --- guided practice on structured
problems using Jupyter notebooks; \textbf{Challenge} (20 min) ---
open-ended problem connecting lecture content to the student's own
SF-PBL project system.

% ============================================================
%  3. METHODOLOGY
% ============================================================
\section{Methodology}

\subsection{Research Design}

This study employed a \textbf{single-cohort pre--post
mixed-methods design} \citep{creswell2018designing} comprising:

\begin{itemize}[noitemsep]
  \item \textit{Quantitative strand}: Pre--post surveys (10-item
    self-efficacy Likert scale, 5-item knowledge MCQ, 10-item PBL
    satisfaction scale), categorical AI usage data, and six-judge
    exhibition rubric scores

  \item \textit{Qualitative strand}: Thematically coded open-ended
    reflections (OE1--OE3) from the post-survey

  \item \textit{Triangulation}: Quantitative learning gains
    corroborated by qualitative breakthrough narratives referencing
    specific SF-PBL phases
\end{itemize}

\subsection{Course Context and Participants}

Course parameters and participant details are presented in
Table~\ref{tab:context}.

\begin{table}[ht]
\centering
\caption{Course Context and Participant Summary}
\label{tab:context}
\begin{tabular}{ll}
\toprule
\textbf{Parameter} & \textbf{Value} \\
\midrule
Course code                    & 702MH0C023 --- Dynamic System Modeling \\
Institution                    & NMIMS MPSTME, Mumbai (NBA-accredited) \\
Programme                      & B.Tech Mechatronics Engineering, Sem.\,IV \\
Academic year                  & 2025--26 \\
Course dates                   & 2 January 2026 -- 25 April 2026 \\
Duration                       & 15 working weeks (3 credits: 2L\,+\,2P/week) \\
N registered                   & 55 students \\
N consenting \& participating  & \textbf{53} (2 did not provide informed consent) \\
N present at exhibition        & \textbf{53} (all consenting participants) \\
Teams                          & 17 (2--4 members per team) \\
Engineering project domains    & 5 \\
\bottomrule
\end{tabular}
\end{table}

Participation in the research survey was voluntary and anonymised;
course grades were not affected by research participation. Informed
consent was obtained from all 53 participants prior to pre-survey
administration. Two registered students did not provide consent
and are excluded from all analyses ($N$\,=\,53 throughout).

\subsection{The SF-PBL Framework: Four-Phase Computational Algorithm}

The SF-PBL project ran concurrently with the 15-week course (2
January -- 25 April 2026) and culminated in a public Mini Project
Exhibition on \textbf{25 April 2026} at NMIMS MPSTME, organised in
collaboration with the NMIMS IUCEE Student Chapter. Each of the
17 teams executed the following four-phase pipeline:

\begin{description}[leftmargin=0pt]
  \item[\textbf{Phase 1: System Identification} (Weeks 1--4).]
    Students selected a real engineering system (student-chosen,
    not prescribed), identified degrees of freedom (DOF) and
    constraints, and defined coordinate frames.
    \textit{Deliverable}: System description + kinematic sketch.
    CO1 | Bloom's Apply (L3).

  \item[\textbf{Phase 2: Mathematical Modelling} (Weeks 5--9).]
    Students drew Free Body Diagrams (FBDs) by hand, applied
    Newton's Second Law or the Lagrangian formulation, derived
    coupled Equations of Motion (EOMs), and converted to
    state-space form $\dot{\mathbf{x}} = f(\mathbf{x},\mathbf{u})$.
    \textit{Deliverable}: Analytical EOMs + state-space
    representation. CO2, CO3, CO4 | Bloom's Analyse--Evaluate
    (L4--L5).

  \item[\textbf{Phase 3: Computational Simulation} (Weeks 10--13).]
    Students implemented the ODE solver in Python using
    \texttt{scipy.integrate.solve\_ivp}, generated time-domain
    plots (position, velocity, energy vs.\ time), and tuned
    control parameters (e.g., PID gains).
    \textit{Deliverable}: Working simulation + validated plots.
    CO5 | Bloom's Create (L6).

  \item[\textbf{Phase 4: Validation \& Communication} (Weeks 14--15).]
    Students validated simulations against energy conservation,
    limit cases, or published benchmarks, then presented to the
    six-judge expert panel at the public exhibition.
    \textit{Deliverable}: Exhibition presentation + validation
    evidence. CO5 | Bloom's Evaluate--Create (L5--L6).
\end{description}

\subsection{The Six-Judge Expert Evaluation Panel}

To ensure authentic, multi-perspective, and bias-reduced
assessment, a six-member external expert panel was constituted
spanning five professional domains (Table~\ref{tab:judges}).

\begin{table}[ht]
\centering
\caption{Six-Judge Expert Evaluation Panel (25 April 2026)}
\label{tab:judges}
\begin{tabular}{clllr}
\toprule
\textbf{\#} & \textbf{Judge} & \textbf{Affiliation} &
\textbf{Domain} & \textbf{Max} \\
\midrule
1 & Kamlesh Panchal         & Manufacturing Industry & Production Engg.    & /16 \\
2 & Hasti Chandarana         & IP Consultancy         & Patent \& Innovation & /16 \\
3 & Nitu Gupta               & External Academic      & Mathematics          & /16 \\
4 & Mahendra Kane (Siemens)  & Siemens Ltd.           & Industry 4.0 / DT    & /20 \\
5 & Parminder Jandoo         & Training Organisation  & Communication        & /16 \\
6 & Debasis Dash             & HR Consultancy         & Behavioural Science  & /16 \\
\midrule
  & \textbf{TOTAL}           &                        &                      & \textbf{/100} \\
\bottomrule
\end{tabular}
\end{table}

\subsection{Student Projects: 17 Teams, 5 Engineering Domains}

The 17 student teams were distributed across five project
categories: (A) Robotics \& Locomotion --- 7 teams (humanoid,
hexapod, multi-terrain robot, quadruped, pick-and-place rover,
robotic gripper, H$_2$ fuel cell UAV); (B) Aerospace \& Marine
--- 3 teams (rocket landing \& docking, ROV, fixed-wing UAV);
(C) Control Systems --- 2 teams (ball balancing Stewart platform,
self-balancing robot); (D) Vehicle Dynamics --- 2 teams (critical
skid velocity, shock absorber); (E) Electromechanical --- 3 teams
(digital twin DC motors, motor modelling, automated EoT crane).

\subsection{Survey Instruments}

\textbf{Self-Efficacy Scale (SE1--SE10):} Ten items rated on a
5-point Likert scale (1\,=\,Strongly Disagree, 5\,=\,Strongly
Agree), each aligned to a specific CO, covering kinematic
analysis, FBD construction, Newton's EOM derivation, Lagrangian
formulation, state-space conversion, Python ODE implementation,
energy validation, PID tuning, plot interpretation, and technical
communication.

\textbf{Knowledge MCQ (PK1--PK5):} Five items assessing degrees
of freedom, Lagrangian formulation, pendulum natural frequency,
viscous damping effects, and Newton's rotation law. Scored 0--5.

\textbf{PBL Experience Scale (PBL1--PBL10):} Ten items on a
5-point scale, grouped into perceived value (PBL1--3), scaffold
quality (PBL4--7), and overall satisfaction (PBL8--10).

\textbf{AI Tool Usage:} Multi-select categorical (ChatGPT, GitHub
Copilot, Claude, None, Other); phase of primary use; estimated
percentage of mathematical derivation completed without AI.

\textbf{Open-Ended Reflections (OE1--OE3):} (OE1) Most valuable
learning; (OE2) Suggested changes for future cohorts; (OE3)
Breakthrough moment when simulation made sense.

\subsection{Data Analysis}

\begin{itemize}[noitemsep]
  \item \textbf{Paired $t$-tests} (two-tailed, $\alpha$\,=\,.05):
    Pre--post SE and PK comparisons

  \item \textbf{Cohen's $d$} (pooled SD:
    $d = \Delta M \,/\, \sqrt{(SD_1^2 + SD_2^2)/2}$):
    small\,=\,0.2, medium\,=\,0.5, large\,=\,0.8
    \citep{cohen1988statistical}

  \item \textbf{Hake's normalised gain}:
    $g = (\text{post}-\text{pre})\,/\,(\text{max}-\text{pre})$ per
    student; bands: low $<$ 0.30, medium 0.30--0.69, high
    $\geq$ 0.70 \citep{hake1998interactive}

  \item \textbf{95\% CI}: $\pm\,1.96 \times SE_{\text{mean}}$

  \item \textbf{Wilcoxon signed-rank test}: non-parametric
    confirmation ($\alpha$\,=\,.05)

  \item \textbf{Post-hoc power analysis} (G*Power framework)

  \item \textbf{Thematic analysis} \citep{thomas2008methods}:
    OE1--OE3 coded by frequency of theme occurrence
\end{itemize}

All analyses were conducted in Python~3.10 using SciPy~1.11 and
NumPy~1.25. Analysis scripts are available from the corresponding
author on request.

% ============================================================
%  4. RESULTS
% ============================================================
\section{Results}

\subsection{Self-Efficacy: Pre--Post Comparison}

Self-efficacy in dynamic system modelling increased significantly
from pre- to post-exhibition (Table~\ref{tab:se_main}). The mean
gain of $+$0.88 points (Cohen's $d$\,=\,2.72) represents a very
large effect substantially exceeding the threshold for practical
significance \citep{cohen1988statistical}. The result was
confirmed by the non-parametric Wilcoxon signed-rank test
($W$\,=\,0.0, $p$\,$<$\,.001) with post-hoc power\,=\,1.000.

\begin{table}[ht]
\centering
\caption{Self-Efficacy Pre--Post Comparison ($N$\,=\,53)}
\label{tab:se_main}
\begin{tabular}{lcccccc}
\toprule
\textbf{Measure} & \textbf{Pre $M$ ($SD$)} & \textbf{Post $M$ ($SD$)} &
\textbf{Gain} & \textbf{95\% CI} & \textbf{$t$(52)} &
\textbf{$p$ \quad $d$} \\
\midrule
SE Mean (1--5) & 2.92 (0.21) & 3.80 (0.40) & $+$0.88 &
[0.75,\,1.00] & 13.91 & $<$.001 \quad 2.72 \\
\bottomrule
\end{tabular}
\end{table}

All ten individual SE items were statistically significant at
$p$\,$<$\,.001 (Table~\ref{tab:se_items}). The largest gains were
observed for SE8 (PID Tuning, $d$\,=\,2.22) and SE3 (Newton's
EOM, $d$\,=\,2.14), both directly activated in Phases 2 and 3.
The smallest gain, SE5 (State-Space, $d$\,=\,1.46), remains a
very large effect and reflects the advanced junction between
analytical derivation and computational implementation.

\begin{table}[ht]
\centering
\caption{Per-Item Self-Efficacy Results ($N$\,=\,53). All $p$\,$<$\,.001.}
\label{tab:se_items}
\begin{tabular}{llrrrr}
\toprule
\textbf{Item} & \textbf{Competency Focus} &
\textbf{Pre $M$} & \textbf{Post $M$} & \textbf{Gain} & \textbf{$d$} \\
\midrule
SE1  & Kinematics              & 2.89 & 3.79 & $+$0.91 & 1.69 \\
SE2  & Free Body Diagram       & 2.94 & 3.77 & $+$0.83 & 1.72 \\
SE3  & Newton's EOM            & 2.89 & 3.81 & $+$0.92 & 2.14 \\
SE4  & Lagrangian              & 2.93 & 3.81 & $+$0.89 & 1.85 \\
SE5  & State-Space             & 2.93 & 3.72 & $+$0.79 & 1.46 \\
SE6  & ODE / Python            & 2.93 & 3.91 & $+$0.98 & 1.82 \\
SE7  & Energy Validation       & 3.02 & 3.77 & $+$0.76 & 1.63 \\
SE8  & PID Tuning              & 2.91 & 3.83 & $+$0.92 & 2.22 \\
SE9  & Plot Interpretation     & 2.89 & 3.77 & $+$0.89 & 1.69 \\
SE10 & Technical Communication & 2.91 & 3.77 & $+$0.87 & 1.57 \\
\bottomrule
\multicolumn{6}{l}{\footnotesize Cohen's $d$ range: 1.46 (SE5)
-- 2.22 (SE8). All $d$ values large by Cohen (1988).}
\end{tabular}
\end{table}

\subsection{Conceptual Knowledge: Pre--Post Comparison}

Knowledge MCQ scores increased significantly (Table~\ref{tab:pk}).
Hake's normalised gain $g$\,=\,0.50 falls in the \textbf{medium}
interactive-engagement category \citep{hake1998interactive}
(0.30\,$\leq g <$\,0.70), in the upper half of the medium band.
The MCQ distribution shift --- from left-skewed (pre, mode\,=\,2,
no student achieving 5/5) to right-skewed (post, 12 students
achieving 5/5) --- confirms broad-based conceptual development
(Table~\ref{tab:pk_dist}).

\begin{table}[ht]
\centering
\caption{Knowledge MCQ Pre--Post Comparison ($N$\,=\,53)}
\label{tab:pk}
\begin{tabular}{lcccccc}
\toprule
\textbf{Measure} & \textbf{Pre $M$ ($SD$)} & \textbf{Post $M$ ($SD$)} &
\textbf{Gain} & \textbf{95\% CI} &
\textbf{$t$(52) \quad $p$} & \textbf{$d$ \quad $g$} \\
\midrule
PK Total (0--5) & 1.96 (0.94) & 3.28 (1.25) & $+$1.32 &
[1.09,\,1.55] & 11.32 \quad $<$.001 & 1.20 \quad 0.50 \\
\bottomrule
\end{tabular}
\end{table}

\begin{table}[ht]
\centering
\caption{MCQ Score Distribution ($N$\,=\,53)}
\label{tab:pk_dist}
\begin{tabular}{lrrrrrr}
\toprule
 & \textbf{$k$=0} & \textbf{$k$=1} & \textbf{$k$=2} &
\textbf{$k$=3} & \textbf{$k$=4} & \textbf{$k$=5} \\
\midrule
Pre  ($n$) & 3 & 11 & 28 & 7 & 4 & 0  \\
Post ($n$) & 0 & 4  & 11 & 16 & 10 & 12 \\
\bottomrule
\end{tabular}
\end{table}

\subsection{PBL Experience Satisfaction}

Post-exhibition PBL satisfaction was consistently high across all
three sub-scales (Table~\ref{tab:pbl}), with a narrow SD of 0.30
indicating high cohort-level consistency in the SF-PBL experience.

\begin{table}[ht]
\centering
\caption{PBL Experience Scale --- Post-Survey ($N$\,=\,53)}
\label{tab:pbl}
\begin{tabular}{llcc}
\toprule
\textbf{Sub-scale} & \textbf{Items} & \textbf{$M$} & \textbf{$SD$} \\
\midrule
PBL Value            & PBL1--3  & 3.70 & ---  \\
Scaffold Quality     & PBL4--7  & 3.76 & ---  \\
Overall Satisfaction & PBL8--10 & 3.76 & ---  \\
\textbf{Total Mean}  & PBL1--10 & \textbf{3.74} & \textbf{0.30} \\
\bottomrule
\end{tabular}
\end{table}

\subsection{Exhibition Rubric Scores: Six-Judge Panel}

The combined six-judge assessment of 17 teams ($N$\,=\,53 present
students) yielded $M$\,=\,82.62/100, $SD$\,=\,9.98,
range\,=\,53--96 (Table~\ref{tab:rubric}).\footnote{One registered
student was absent on exhibition day and scored 0/100; this
student is excluded from all combined statistics. Including the
absent student: $M$\,=\,81.09, $SD$\,=\,14.97, $N$\,=\,54.} The
SD of approximately 10 points reflects meaningful discriminative
ability while confirming that all present students achieved a
substantive standard.

\begin{table}[ht]
\centering
\caption{Six-Judge Panel Scores ($N$\,=\,53 students, 17 teams,
  25 April 2026)}
\label{tab:rubric}
\begin{tabular}{lllrrr}
\toprule
\textbf{Judge} & \textbf{Domain} & \textbf{Max} &
\textbf{$M$} & \textbf{$SD$} & \textbf{Range} \\
\midrule
Kamlesh Panchal         & Manufacturing          & /16 & 13.94 & 1.23 & 11--16 \\
Hasti Chandarana        & IP / Patent            & /16 & 12.53 & 2.68 & 8--16  \\
Nitu Gupta              & Mathematics (External) & /16 & 12.77 & 2.67 & 6--16  \\
Mahendra Kane (Siemens) & Industry 4.0           & /20 & 17.09 & 2.18 & 12--20 \\
Parminder Jandoo        & Communication          & /16 & 13.26 & 2.13 & 5--16  \\
Debasis Dash            & Behavioural Science    & /16 & 13.02 & 2.28 & 8--16  \\
\midrule
\textbf{COMBINED TOTAL} & & \textbf{/100} &
\textbf{82.62} & \textbf{9.98} & \textbf{53--96} \\
\bottomrule
\end{tabular}
\end{table}

\subsection{AI Tool Usage}

Forty-five of 53 students (84.9\%) reported using AI tools during
the SF-PBL project (Table~\ref{tab:ai}). The most common use case
was Python simulation coding (Phase 3) and debugging; 39.6\%
estimated 50--75\% of their mathematical derivation was completed
without AI assistance, indicating human-led mathematical reasoning
remained dominant throughout.

\begin{table}[ht]
\centering
\caption{AI Tool Usage Distribution ($N$\,=\,53)}
\label{tab:ai}
\begin{tabular}{lrr}
\toprule
\textbf{Tool / Combination} & \textbf{$n$} & \textbf{\%} \\
\midrule
ChatGPT                  & 18 & 34.0 \\
ChatGPT + GitHub Copilot & 11 & 20.8 \\
None                     &  8 & 15.1 \\
Claude                   &  6 & 11.3 \\
ChatGPT + Claude         &  5 &  9.4 \\
Other                    &  5 &  9.4 \\
\midrule
\textbf{Total AI users}  & \textbf{45} & \textbf{84.9} \\
\bottomrule
\end{tabular}
\end{table}

\subsection{Qualitative Themes}

Thematic coding of open-ended post-survey responses revealed
consistent patterns across all three open-ended items ($N$\,=\,53):

\textbf{OE1 --- Most Valuable Learning:} Mathematical validation
of simulation results ($n$\,=\,18, 34.0\%); team collaboration
($n$\,=\,10, 18.9\%); debugging complex models ($n$\,=\,8,
15.1\%); other ($n$\,=\,17, 32.1\%).

\textbf{OE3 --- Breakthrough Moment:} Linking mathematical
derivation to simulation output ($n$\,=\,20, 37.7\%;
e.g., \textit{``The moment my simulation plot matched my
analytical solution, I finally understood what the equations meant
physically''}); understanding PID tuning from first principles
($n$\,=\,12, 22.6\%); seeing the physical model work in simulation
($n$\,=\,11, 20.8\%); other ($n$\,=\,10, 18.9\%).

\textbf{OE2 --- Suggested Changes:} Smaller group sizes of 2--3
members ($n$\,=\,16, 30.2\%); less documentation, more coding time
($n$\,=\,12, 22.6\%); more project development time ($n$\,=\,11,
20.8\%); other ($n$\,=\,14, 26.4\%).

% ============================================================
%  5. DISCUSSION
% ============================================================
\section{Discussion}

\subsection{Interpreting the Self-Efficacy Effect Size}

The self-efficacy gains ($d$\,=\,2.72) substantially exceed those
reported in comparable PBL studies. \citet{freeman2014active}'s
meta-analysis of active learning found a mean $d$\,$\approx$\,0.47;
\citet{dochy2003effects}'s PBL meta-analysis reported
$d$\,$\approx$\,0.65. The very large Cohen's $d$ here warrants
careful interpretation: the narrow pre-survey SD (0.21 on a
5-point scale) reflects a relatively homogeneous cohort at
baseline, which mathematically amplifies the standardised effect
size. The absolute gain ($+$0.88 points, 95\% CI [0.75, 1.00]) is
arguably the more practically meaningful metric, representing a
reliable cohort-wide shift from approximately ``Neutral'' to
``Agree'' across all ten competency items.

Critically, the two items with the largest gains --- SE8 (PID
Tuning, $d$\,=\,2.22) and SE3 (Newton's EOM, $d$\,=\,2.14) ---
correspond directly to Phases 2 and 3 of the SF-PBL pipeline,
confirming that the computational integration mechanism activated
the intended learning outcomes rather than producing uniform,
undifferentiated gains.

\subsection{Knowledge Gains and Hake's Normalised Gain}

Hake's normalised gain $g$\,=\,0.50 places this course in the
\textbf{medium} interactive-engagement category
\citep{hake1998interactive} (0.30\,$\leq g <$\,0.70), in the
upper half of the medium band. This result is particularly notable
given that: (a) the MCQ measured conceptual understanding rather
than procedural recall; (b) students began with low prior knowledge
(pre: $M$\,=\,1.96/5, 39\% correct); and (c) assessment under
examination conditions reduced social desirability effects.

\subsection{AI Integration: Complementarity Not Substitution}

The finding that 84.9\% of students used AI tools --- primarily
for Python simulation coding and debugging --- suggests an emerging
pattern of \textbf{human--AI complementarity} in SF-PBL. The
framework's structural design encouraged this: the rubric
explicitly required hand-drawn FBDs and analytical EOM derivations
\textit{before} simulation (Phase 3), creating a curricular
incentive to use AI as an accelerator for implementation rather
than a substitute for mathematical reasoning. This contributes to
the emerging discourse on AI in STEM education \citep{tlili2023chatgpt}:
SF-PBL channels AI use productively by positioning derivation as
the primary intellectual activity.

\subsection{The Six-Judge Panel as Authentic Assessment}

The six-judge panel spanning five professional domains represents
a significant departure from conventional faculty-only assessment.
The panel mean of 82.62/100 ($SD$\,=\,9.98, range\,=\,53--96)
demonstrates that all present students achieved substantive scores
with meaningful spread across team quality levels, consistent with
\citeauthor{gulikers2004five}'s \citeyearpar{gulikers2004five}
criteria for authentic assessment validity. The Industry 4.0 judge
(Siemens: $M$\,=\,17.09/20\,=\,85.5\%) awarded the highest
relative scores, suggesting student projects met contemporary
digital engineering standards relevant for mechatronics graduates
entering Industry 4.0 workplaces.

\subsection{The Breakthrough Theme: Derivation-to-Simulation}

The dominant qualitative breakthrough theme --- \textit{linking
mathematical derivation to simulation output} ($n$\,=\,20, 37.7\%)
--- directly validates the theoretical premise of SF-PBL: that
computation serves as the bridge between abstract mathematical
formalism and physical intuition. This theme appeared independently
of AI tool usage pattern, project complexity level, and team size,
suggesting that the derivation-to-simulation moment is a
\textbf{universal learning landmark} in computationally-integrated
dynamics education requiring specific instructional scaffolding.

\subsection{Practical Implications for STEM Educators}

The SF-PBL framework offers three directly replicable elements:
(1) the \textbf{four-phase pipeline} is domain-agnostic and
transferable to thermodynamics, fluid mechanics, structural
analysis, or control systems with minor adaptation; (2) the
\textbf{E-T-C lecture model} provides an active-learning structure
within existing contact hours without requiring curriculum
restructuring; (3) the \textbf{six-judge rubric with
CO$\rightarrow$PO traceability} provides an
accreditation-compliant (NBA/ABET) authentic assessment instrument
with inter-professional credibility that faculty-only panels cannot
achieve.

\subsection{Limitations}

Several limitations should be acknowledged. First, the study
employed a single-group pre--post design without a concurrent
control condition; causal attribution to the SF-PBL framework
versus general course maturation effects is therefore constrained.
Second, the sample ($N$\,=\,53) is from a single institution,
programme, and cohort, limiting generalisability. Third, AI usage
data is self-reported and subject to social desirability bias.
Fourth, formal inter-rater reliability (Cohen's $\kappa$) between
the six expert judges was not computed, representing a measurement
limitation for the rubric data. Fifth, the narrow pre-survey SD
suggests possible range restriction in the self-efficacy baseline.
Future research should address these limitations through controlled
comparisons, multi-institution replication, longitudinal follow-up,
and formal inter-rater reliability analysis for expert panels.

% ============================================================
%  6. CONCLUSION
% ============================================================
\section{Conclusion}

This study presented the design, implementation, and mixed-methods
evaluation of the SF-PBL framework in an undergraduate engineering
dynamics course at an NBA-accredited institution. The framework
operationalises a four-phase computational pipeline --- System
Identification, Mathematical Modelling, Python Simulation,
Validation \& Communication --- within a 15-week OBE-aligned
course assessed by a six-judge external expert panel.

Principal findings:
\begin{itemize}[noitemsep]
  \item \textbf{Very large self-efficacy gains} across all ten
    competency items ($d$\,=\,2.72, $p$\,$<$\,.001), largest for
    PID Tuning ($d$\,=\,2.22) and Newton's EOM ($d$\,=\,2.14)

  \item \textbf{Large knowledge gains} with medium Hake normalised
    gain ($g$\,=\,0.50, upper half of medium band; $d$\,=\,1.20,
    $p$\,$<$\,.001)

  \item \textbf{High PBL satisfaction} ($M$\,=\,3.74/5,
    $SD$\,=\,0.30) with high cohort consistency

  \item \textbf{Strong exhibition performance}
    ($M$\,=\,82.62/100, $SD$\,=\,9.98) across 17 project teams
    in five engineering domains

  \item \textbf{Emergent human--AI complementarity}: 84.9\% AI
    usage concentrated in simulation coding; mathematical
    derivation predominantly human-led

  \item \textbf{Universal breakthrough moment}: Linking analytical
    EOM derivation to simulation output --- the most frequently
    reported transformative learning experience (37.7\%)
\end{itemize}

The SF-PBL framework provides a replicable, evidence-based, and
OBE-compliant model for STEM educators seeking to integrate
computational simulation authentically into engineering dynamics
pedagogy. All instruments, rubrics, and analysis scripts are
available from the corresponding author upon reasonable request.

% ============================================================
%  STATEMENTS
% ============================================================
\section*{Data Availability Statement}
The anonymised individual-level survey dataset ($N$\,=\,53, 65 variables),
exhibition rubric evaluation scores (six judges across 17 teams), survey instruments,
rubric specifications, Python analysis scripts (SciPy~1.11, NumPy~1.25), and
representative student project artefacts are publicly available in the official
repository: \url{https://github.com/sunny-nanade/SF-PBL-Engineering-Dynamics-2026}.

\section*{Ethics Statement}

Informed consent was obtained from all participants prior to
pre-survey administration. Participation was voluntary; course
grades were not affected by non-participation. Data were
anonymised for analysis. The study was conducted in accordance
with institutional guidelines for educational research at NMIMS
MPSTME.

\section*{Author Contributions}

SN: Conceptualisation, course design, SF-PBL framework
development, data collection, formal analysis, methodology,
writing (original draft, review \& editing), project
administration.

\section*{Conflict of Interest}

The author declares no conflict of interest.

\section*{Acknowledgements}

The author acknowledges the contribution of the six expert judges,
the NMIMS IUCEE Student Chapter (exhibition organisation), and the
53 students of DSM Batch 2025--26 whose project work made this
research possible.

% ============================================================
%  REFERENCES
% ============================================================
\bibliographystyle{apalike}
\bibliography{DSM_references}

% ============================================================
%  APPENDIX
% ============================================================
\appendix
\section{Survey Instruments}
\label{app:instruments}

\subsection*{A.1 Self-Efficacy Items (SE1--SE10)}

\textit{Instructions: Rate your confidence in each ability on a
5-point scale (1\,=\,Strongly Disagree, 5\,=\,Strongly Agree).}

\begin{enumerate}[label=SE\arabic*.]
  \item I can apply kinematic analysis to identify DOF in a
    multi-body system. [CO1]
  \item I can construct a correct Free Body Diagram for a complex
    mechanism. [CO2]
  \item I can derive equations of motion using Newton's Second
    Law. [CO3]
  \item I can formulate the Lagrangian for a conservative dynamic
    system. [CO4]
  \item I can convert a second-order ODE system to state-space
    form. [CO5]
  \item I can implement a numerical ODE solver
    (\texttt{solve\_ivp}) in Python. [CO5]
  \item I can validate my simulation using energy conservation
    checks. [CO5]
  \item I can tune PID controller gains to achieve a desired
    dynamic response. [CO5]
  \item I can interpret simulation time-domain plots of position,
    velocity, and energy. [CO5]
  \item I can communicate my simulation results clearly to a
    non-specialist. [Communication]
\end{enumerate}

\subsection*{A.2 Knowledge MCQ Items (PK1--PK5)}

\textit{Instructions: Select the single best answer.}

\begin{enumerate}[label=PK\arabic*.]
  \item A rigid body constrained to rotate in a plane has:
    (a) 3 DOF \quad (b) \textbf{1 DOF} \quad (c) 2 DOF \quad
    (d) 0 DOF

  \item The Lagrangian $L$ is defined as:
    (a) $L = V - T$ \quad \textbf{(b) $L = T - V$} \quad
    (c) $L = T + V$ \quad (d) $L = \partial T/\partial q$

  \item The natural frequency of a simple pendulum (length $l$):
    \textbf{(a) $\omega_n = \sqrt{g/l}$} \quad
    (b) $\omega_n = \sqrt{l/g}$ \quad
    (c) $\omega_n = g/l$ \quad (d) $\omega_n = l/g$

  \item Adding viscous damping to an oscillator:
    (a) Increases natural frequency \quad
    (b) Has no effect \quad
    \textbf{(c) Decreases the oscillation frequency} \quad
    (d) Doubles the amplitude

  \item Newton's Second Law for rotation states:
    (a) $\tau = ma$ \quad \textbf{(b) $\tau = I\alpha$} \quad
    (c) $\tau = I\omega$ \quad (d) $\tau = m\alpha$
\end{enumerate}

\end{document}
